%% -----------------------------------------------------------
\section{Risks Facing Marginalized Communities}
%% -----------------------------------------------------------

\subsection{Accessibility Exclusion}

Accessibility exclusion constitutes the most structurally entrenched risk. Platforms that ostensibly expand care actively block disabled users from accessing it. A study with 93 blind users documented systematic accessibility failures across \gls{dmh} platforms despite legal requirements and available technical solutions~\cite{khan_seo_2025}. Most \gls{dmh} tracking services lacked screen reader compatibility, used visual-only displays without text alternatives, and created interaction patterns impossible to navigate without sight. Participants described forced choices between inaccessible mainstream applications with full features or specialized disability-focused platforms with limited functionality~\cite{khan_seo_2025}. This affects a population that already experiences depression and anxiety at two to three times the rates of sighted individuals.

The pattern mirrors decades of web accessibility failure: legal requirements exist, implementation is ignored, and disabled populations are structurally locked out of technologies claimed to democratize access~\cite{khan_seo_2025}.

\subsection{Demographic Bias in Large Language Models}

\gls{llm} demographic bias is empirically documented. Omar et al.~\cite{omar_2025} conducted a systematic review of 24 studies evaluating demographic biases in medical \gls{llm}s and found that 22 (91.7\%) identified biases, with gender bias reported in 15 of 16 studies (93.7\%) and racial or ethnic biases in 10 of 11 studies (90.9\%). Pfohl et al.~\cite{pfohl_2024} developed a multifactorial framework for assessing equity-related \gls{llm} harms and conducted a large-scale empirical study with the Med-PaLM 2 \gls{llm}, involving 806 raters from diverse demographic backgrounds and generating over 17,000 human ratings; their findings revealed systematic bias patterns that narrower evaluation approaches missed. Lawrence et al.~\cite{lawrence_2024} identified additional dimensions of \gls{llm} risk in mental health contexts, including the absence of standardized evaluation protocols and compounding effects when \gls{llm}s are deployed at scale.

These findings carry particular weight for mental health applications. In clinical presentations, race, gender, sexuality, disability status, and socioeconomic status are directly implicated in symptom expression, help-seeking behavior, treatment response, and likelihood of accurate diagnosis. Bias on these dimensions in mental health AI carries direct clinical consequences~\cite{timmons_2023,haider_2024}.

Straw and Callison-Burch~\cite{straw_2020} specifically examined language model bias in mental health contexts, demonstrating that general-purpose \gls{llm}s reproduce racial and gender stereotypes in mental health-relevant outputs, providing a mechanistic account of how bias encoded in training data propagates into care-relevant advice. The broader review by Haider et al.~\cite{haider_2024} corroborates this at scale across 30 studies, finding consistent racial disparities in AI-driven healthcare recommendations.

\subsection{Neurotypical Bias Against Neurodivergent Users}

A distinct category of bias affects neurodivergent users: neurotypical bias, in which \gls{llm}s trained predominantly on neurotypical-authored text encode assumptions about emotional processing, social communication, and help-seeking that may conflict with neurodivergent experience~\cite{killian_2023}. Responses considered empathic and clear by neurotypical raters were judged by clinicians working with autistic populations as ``overly wordy, vague, and overwhelming,'' reflecting a mismatch between assumed and actual cognitive needs~\cite{killian_2023}. For a population already prone to misunderstanding in clinical encounters, AI systems that replicate these misunderstandings at scale represent a structural barrier to care~\cite{papadopoulos_2025}.

Additionally, the risk of emotional over-reliance is particularly pronounced for autistic users. Papadopoulos~\cite{papadopoulos_2025} documents cases in which autistic individuals developed primary social bonds with AI companions, to the detriment of human connection, and in one documented case a companion AI failed to intervene during acute suicidal ideation, reflecting the same crisis detection failure documented across populations by De Freitas et al.~\cite{de_freitas_2024}.

\subsection{Online Safety and Harassment}

Online safety risks disproportionately burden \gls{lgbtq} youth using digital platforms for mental health support. An analysis of 173 youth Instagram users (86 \gls{lgbtq}, 87 heterosexual) revealed that \gls{lgbtq} youth experienced significantly more high-risk online interactions: 68\% experienced sexual messages from strangers compared to 23\% of heterosexual youth ($p < 0.001$), and 34\% experienced online interactions directly linked to self-harm behaviors compared to 11\% of heterosexual youth ($p < 0.001$)~\cite{tanni_2024}. Each additional high-risk interaction category correlated with 1.4-point increases in PHQ-9 depression scores and 1.8-point increases in \gls{gad} anxiety scores~\cite{tanni_2024}. Existing AI content moderation demonstrated dual failures: only 23\% of reported harassment resulted in content removal, while 41\% of identity-affirming posts were flagged as inappropriate~\cite{tanni_2024}.

\subsection{Crisis Mismanagement and Harmful Advice}

Crisis recognition and response represent the most acutely dangerous risk domain. Mental health crises appeared in 17\% of conversations with companion AIs ($n = 289$), including suicidal ideation (8\%), severe depression (12\%), and self-harm discussion (9\%)~\cite{de_freitas_2024}. Performance testing found chatbots failed to recognize crisis situations in 62\% of test scenarios involving suicidal ideation and provided appropriate escalation responses in only 31\% of detected crises~\cite{de_freitas_2024}. In 23\% of crisis conversations, rigid safety guardrails rejected users during moments of vulnerability with generic responses, potentially increasing distress rather than mitigating it~\cite{de_freitas_2024}.

The real-world stakes of crisis mismanagement are illustrated by the case of \gls{neda}'s Tessa chatbot. Deployed as a peer support tool, Tessa was shut down in May 2023 after users reported it provided harmful weight-loss advice, recommending calorie restriction and weight tracking in direct contradiction to evidence-based eating disorder treatment~\cite{neda_2023}. The case exemplifies De Freitas et al.'s empirical findings: AI trained on general data lacks the clinical knowledge to navigate complex mental health situations safely, and the assumption that accessibility automatically benefits vulnerable users is demonstrably false~\cite{de_freitas_2024}.

Palmer et al.~\cite{palmer_2025} add a regulatory dimension, documenting that none of the three dominant regulatory models for digital mental health tools (laissez-faire, highly regulated, and the current \gls{fda} approach) satisfactorily balances the promise of access with ethical risk management, leaving a structural gap through which harmful deployments continue to reach vulnerable populations~\cite{palmer_2025}.

\subsection{Ethical Violations in Therapeutic AI Deployment}

Meadi et al.~\cite{meadi_2025} systematically examined the ethical challenges of deploying conversational AI ``as therapist.'' Their analysis identified violations across autonomy, beneficence, non-maleficence, and justice, finding that platforms routinely overclaim therapeutic capability, inadequately disclose limitations, and fail to provide meaningful alternatives when users require escalated care. Marginalized populations bear elevated exposure to these ethical failures because they have fewer alternatives when AI interventions fail~\cite{meadi_2025}.

\subsection{Implementation Gaps and Sustainability Risks}

Implementation gaps create systemic vulnerabilities. An expert panel review identified that current AI mental health tools lack standardized safety testing protocols, real-time monitoring systems, clear crisis escalation pathways, long-term outcome tracking, and mechanisms to identify when users discontinue due to harm~\cite{smith_2023}. Dehbozorgi et al.~\cite{dehbozorgi_2025} corroborate these findings, establishing that ethical frameworks and safety protocols remain underdeveloped relative to deployment scale. A scoping review additionally identified sustainability crises, with chatbot companies lacking viable business models: when platforms shut down, vulnerable users face abrupt service termination with no continuity of care and no human alternative~\cite{casu_2024}.
